La Figura 1 muestra el comportamiento de la temperatura del sistema en función del tiempo. Se observa una disminución continua de la temperatura, con una pendiente inicialmente más pronunciada y luego cada vez menor. Este comportamiento es consistente con la ley de enfriamiento de Newton, ya que la rapidez de cambio de la temperatura depende de la diferencia entre la temperatura del cuerpo y la del ambiente. Al inicio, cuando esta diferencia es mayor, el enfriamiento es más rápido; posteriormente, al acercarse al equilibrio térmico, la variación se hace más lenta.

Por otra parte, la Figura 2 presenta el logaritmo de la temperatura relativa \(\ln(T-T_a)\) en función del tiempo. En este caso se obtiene una relación aproximadamente lineal, lo cual confirma que el proceso sigue un comportamiento exponencial. La linealización de los datos permite verificar experimentalmente el modelo teórico y facilita la obtención de la constante de enfriamiento a partir de la pendiente de la recta.

\subsection{Análisis}

Los resultados indican que el sistema intercambia energía térmica con el ambiente hasta tender al equilibrio. El hecho de que la temperatura descienda de forma exponencial refleja que la tasa de pérdida de calor no es constante, sino proporcional a la diferencia térmica con el entorno. Esto concuerda con el modelo de Newton, en el cual el cuerpo pierde calor más rápidamente cuando su temperatura inicial es mucho mayor que la del medio.

La linealidad observada en la gráfica logarítmica sugiere que, dentro del intervalo de tiempo analizado, el sistema se comportó de manera cercana a las condiciones ideales del modelo. En particular, se puede considerar que el ambiente permaneció aproximadamente constante y que los mecanismos de transferencia de calor estuvieron dominados por el intercambio con el entorno.

Los datos experimentales muestran un enfriamiento compatible con la ley de Newton: una disminución exponencial de la temperatura y una linealidad clara en la representación logarítmica. El alto valor de \(r^2\) respalda fuertemente el modelo, por lo que se puede concluir que la ley describe de forma satisfactoria el comportamiento térmico observado en el experimento.

\subsection{Incertidumbres y fuentes de error}

Aunque el ajuste obtenido es excelente, existen varias fuentes de incertidumbre que pueden afectar los resultados. En primer lugar, la precisión del instrumento de medición de temperatura limita la exactitud de cada dato experimental. Además, pequeñas variaciones en la temperatura ambiente pueden modificar la diferencia térmica y alterar el proceso de enfriamiento.

También pueden influir errores en el tiempo de registro, la ubicación del sensor, el contacto incompleto entre el termómetro y el sistema, o una posible distribución no uniforme de la temperatura dentro del cuerpo. Otro factor importante es que el modelo supone un ambiente ideal y una transferencia de calor uniforme, condiciones que en la práctica no siempre se cumplen totalmente.

En conjunto, estas incertidumbres explican las pequeñas desviaciones respecto al comportamiento ideal, aunque no afectan de manera significativa la tendencia general de los resultados.

